\section{Seer-v2 and the interventions}\label{sec:method}

\subsection{A fixed restoration backbone}\label{sec:method:pixel}

The pixel branch is a five-scale Restormer-style U-Net~\citep{restormer}
with $22.4$\,M trainable parameters. Its widths are
$(20,40,80,160,320)$, encoder block counts $(2,4,4,6,8)$, and decoder block
counts $(6,4,4,2)$. Frame-wise blocks perform most of the reconstruction;
a residual 3-D convolution and cross-frame attention at the coarsest scale
provide temporal mixing. Local illumination cues---luma, channel maximum,
local mean, low-pass luma, and gradient magnitude---are projected into each
encoder and decoder scale.
A zero-initialised residual head gives
$\Yhat=\X+\mathcal P_\theta(\X;Z)$.

\begin{figure}[t]
\centering
\resizebox{\columnwidth}{!}{\begin{tikzpicture}[
  box/.style={draw=linkink,rounded corners=2pt,align=center,
              minimum height=9mm,text width=29mm,font=\small},
  memory/.style={box,fill=blue!5},
  train/.style={box,fill=orange!8},
  arr/.style={->,thick,draw=linkink},
  note/.style={font=\scriptsize,align=center}
]
\node[box] (input) at (0,0) {Low-light clip\\$16\times H\times W$};
\node[memory] (encoder) at (4,1.4) {Frozen V-JEPA\\ViT-L};
\node[memory] (tokens) at (8,1.4) {Blocks\\$5,11,17,23$};
\node[train] (pixel) at (4,-1.1) {Pixel branch\\$22.4$\,M};
\node[box] (output) at (8,-1.1) {Enhanced clip\\$\Yhat=\X+\mathrm{residual}$};
\node[note] (inject) at (6,-0.05) {Fixed depth-to-scale map\\zero-init cross-attention};
\draw[arr] (input) -- (encoder);
\draw[arr] (encoder) -- (tokens);
\draw[arr] (input) -- (pixel);
\draw[arr] (tokens.south) -- ++(0,-0.45) -| (pixel.north);
\draw[arr] (pixel) -- (output);
\node[note] at (4,-2.2) {Training: reconstruction + Retinex auxiliary heads\\
Inference probes: averaging / centred PCA of the token memory};
\end{tikzpicture}
}
\caption{The studied Seer-v2 model. A frozen video encoder supplies four token
fields to a small trainable pixel branch through fixed scale assignments.
Token interventions change this memory while leaving the trained branch fixed.
Auxiliary reflectance and illumination heads are used only during training.}
\label{fig:arch}
\end{figure}

\paragraph{Frozen memory.}
We use the released V-JEPA\,2.1 ViT-L checkpoint distilled from ViT-G
\citep{vjepa2}. The encoder is frozen throughout restoration training.
It receives $16$ frames resized to $256\times256$, producing
$8\times16\times16=2048$ tokens of width $1024$ at each of blocks
$\Lset=\{5,11,17,23\}$ (zero-indexed).
At spatial scales $1/2,1/4,1/8,1/16$, the pixel branch reads these four
fields respectively, with the same assignment in the decoder.
There is no learned mixture across depths.

For frame $t$, $Z_{s,t}$ contains the $256$ spatial tokens from tubelet
$\lfloor t/2\rfloor$ at the assigned depth. Encoder and decoder features at
scale $s$ read this memory through residual multi-head cross-attention,
\begin{equation}\label{eq:xattn}
F_{s,t}\leftarrow F_{s,t}+W^s_o
  \operatorname{Attn}(F_{s,t}W^s_Q,Z_{s,t}W^s_K,Z_{s,t}W^s_V).
\end{equation}
The output projection $W^s_o$ starts at zero. The default attention has no
added query--key positional correspondence; spatial grounding is a separate
intervention.

\paragraph{Objective.}
Training-only heads predict reflectance $\hat R$ and illumination $\hat L$.
The default loss is
\begin{equation}\label{eq:loss}
\mathcal L_0=
1.5\,\rho(\Yhat,\Y)+
0.15\,\rho(\hat R\odot\hat L,\Y)+
0.02\,\mathrm{TV}_{xy}(\hat L),
\end{equation}
where $\rho$ is the mean Charbonnier loss with $\epsilon=10^{-3}$ and
$\mathrm{TV}_{xy}$ is spatial total variation. Inference uses the clipped
residual prediction without exposure or colour calibration.

\subsection{Intervening on the memory}\label{sec:method:interventions}

Spatial averaging replaces each tubelet's tokens with their spatial mean;
temporal averaging replaces each spatial position with its mean over
tubelets. Both broadcast the result to the original grid, preserving the
attention tensor shapes.

For one clip and tap, write $Z\in\mathbb R^{2048\times1024}$ and let
$\mu_Z$ be its mean token. If $V_{Z,k}$ contains the top $k$ principal
directions of $Z-\mathbf1\mu_Z$, the per-clip intervention is
\begin{equation}\label{eq:cliprank}
\mathcal T^{\mathrm{clip}}_k(Z)=
\mathbf1\mu_Z+(Z-\mathbf1\mu_Z)V_{Z,k}V_{Z,k}^{\top}.
\end{equation}
This preserves a full $1024$-dimensional clip mean; only the centred token
variation is restricted to rank $k$.
The fixed intervention uses a mean $\mu_{\rm tr}$ and directions
$V_{{\rm tr},k}$ estimated separately for each tap from training clips,
\begin{equation}\label{eq:globalrank}
\mathcal T^{\mathrm{fixed}}_k(Z)=
\mathbf1\mu_{\rm tr}+
(Z-\mathbf1\mu_{\rm tr})V_{{\rm tr},k}V_{{\rm tr},k}^{\top}.
\end{equation}
The comparison changes both the mean and the basis; it does not isolate basis
adaptivity alone.
An exploratory validation control crosses the two means and the two bases
using the same reconstruction formula, alongside identity and broadcast
clip-mean controls (App.~\ref{app:meanbasis}).
The adaptive directions always come from clip-centred tokens, including when
the reconstruction uses the training-set mean.

\subsection{Testing a temporal objective}\label{sec:method:temporal}

Let $\Delta Z_t=Z_{t+1}-Z_t$ be the block-23 token difference along tubelets.
We add $0.3\mathcal L_{\Delta}$ to Eq.~\eqref{eq:loss}, with gradients through
the frozen encoder to the predicted pixels.
The relative form is
$\mathcal L_{\Delta}^{\mathrm{rel}}=
\|\Delta\Phi(\Yhat)-\Delta\Phi(\Y)\|_1/
\max(\|\Delta\Phi(\Y)\|_1,\epsilon)$.
A temporally constant prediction scores approximately one, so shrinking the
variation can reduce this objective without reconstructing the dynamics.
The cosine form instead uses
\begin{equation}\label{eq:tdcos}
\begin{aligned}
\mathcal L_{\Delta}^{\mathrm{cos}}={}&
1-\overline{\cos(\Delta\Phi(\Yhat),\Delta\Phi(\Y))}\\
&+\left|\log \frac{\max(\|\Delta\Phi(\Yhat)\|_1,\epsilon)}
                       {\max(\|\Delta\Phi(\Y)\|_1,\epsilon)}\right|,
\end{aligned}
\end{equation}
where cosine is computed per token over channels and the $\ell_1$ quantities
are mean absolute magnitudes. We record both directional
agreement and the magnitude ratio to check what the loss actually changes.

\subsection{Training}\label{sec:method:train}

All reported converged arms use a fixed $20{,}000$-step schedule.
The pixel branch starts from scratch with AdamW, learning rate
$2\times10^{-5}$, weight decay $10^{-4}$, betas $(0.9,0.99)$,
batch one and eight-step gradient accumulation.
Training uses random $256\times256$ crops of $16$-frame clips, flips,
gradient clipping at one, bf16 autocast, and gradient checkpointing.
Validation is monitored without selecting a checkpoint; the final checkpoint
is scored. The default objective contains no latent temporal term.
